% !PRINTED 506
\weberpage{506}
%<<CONT
\[
a_1E_1\Phi_1 = P_1, \quad a_2E_2\Phi_2 = P_2
\]
primitive Functionen werden, und dann ist auch
\weq{2}{a_1a_2E_1E_2\Phi = P_1P_2}
eine primitive Function. Aus (1) und (2) folgt
\weq{3}{a_1a_2E_1E_2P = aEP_1P_2}
und mithin
\weq{4}{a = a_1a_2.}
Wenn also $a=1$ ist, so müssen die natürlichen Zahlen $a_1, a_2$ auch $=1$ sein, und wir haben den Satz:

4. \gesperrt{Ist $\omega$ ein ganzes Functional in $\Omega$ und $\Phi(t)$ eine Function von $t$ mit ganzen Coëfficienten in $\overline{R}$, die für $t=\omega$ verschwindet, so hat auch jeder rationale Theiler von $\Phi(t)$ ganze Coëfficienten in $\overline{R}$; insbesondere hat die irreducible Function $F(t)$, von der $\omega$ nach §.\,136, 8. eine Wurzel ist, ganze rationale Functionale zu Coëfficienten.}

Dieser Satz ist deshalb von principieller Bedeutung, weil er lehrt, dass ein Functional $\omega$, das im Körper $\Omega$ zu den gebrochenen gehört, nicht etwa in einem höheren Körper als ganzes Functional erscheinen kann.

Wenn ein Functional $\omega$ nicht ganz ist, so genügt es einer Gleichung
\[
\Phi(\omega) = \omega^m + A_1\omega^{m-1} + A_2\omega^{m-2} + \cdots + A_m = 0,
\]
worin die Coëfficienten $A_1, A_2, \ldots, A_m$ zwar rationale, aber nicht ganze Functionale sind. Ist $a$ der Hauptnenner der absoluten Werthe von $A_1, A_2, \ldots, A_m$, so sind $aA_1, aA_2, \ldots, aA_m$ ganze Functionale. Nun ist
\[
a^m\Phi(\omega) = (a\omega)^m + aA_1(a\omega)^{m-1} + a^2A_2(a\omega)^{m-2} + \cdots
+ a^mA_m = 0,
\]
und $a\omega$ ist daher ein ganzes Functional. Daraus folgt:

5. \gesperrt{Jedes Functional $\omega$ des Körpers $\Omega$ lässt sich durch Multiplication mit einer ganzen rationalen Zahl in ein ganzes Functional verwandeln, und daher kann man jedes Functional $\omega$ als Quotienten zweier ganzer Functionale darstellen. Diese Darstellung ist auf unendlich viele verschiedene Arten möglich, unter anderem so, dass der Nenner eine natürliche Zahl ist.}
